\chapter{关键设计与实现特点}

前面章节已经说明 MediFlow 的背景、整体设计、三项任务、实现和测试结果。
本章从系统设计角度提炼五个贯穿三项任务的技术决策，
说明它们分别解决了什么问题、在本项目中如何落地、
以及给整个系统带来了哪些具体收益。

\section{MCP 工具与平台边界}

赛题要求基于 Nexent 构建智能体，基于 DataMate 完成数据处理。
项目实现将平台访问、领域规则和智能体配置拆成三层：
Nexent 只负责理解用户任务和选择工具，
不接触 DataMate 的接口细节；
MCP 工具服务负责参数校验和结果整理，
内部再调用 DataMate 客户端和领域处理模块。
三个已发布智能体和自定义智能体均通过 MCP 工具名称绑定相应功能，
工具的参数约定和服务端校验保持一致。

\begin{table}[H]
\centering
\caption{三层职责划分}
\label{tab:three-layer}
\begin{tabularx}{\textwidth}{L{2.8cm}L{4.0cm}Y}
\toprule
\textbf{层次} & \textbf{运行载体} & \textbf{负责的内容} \\
\midrule
任务理解与对话 & Nexent 智能体 & 判断用户意图、选择工具、组织回复 \\
\addlinespace
协议与编排 & MCP 工具服务 & 参数检查、流程编排、结果汇总 \\
\addlinespace
平台执行与领域规则 &
DataMate + MediFlow 服务 &
数据集操作、清洗任务、知识抽取、SQL 安全 \\
\bottomrule
\end{tabularx}
\end{table}

因此，DataMate 接口变化由客户端适配层承接，
领域规则变化由服务模块承接，
智能体配置只保存角色提示、约束提示和工具绑定。
新增智能体时绑定已有 MCP 工具即可，
无需在提示词中重复写入 DataMate 调用细节。

\section{阶段标识与结果追踪}

数据处理、知识构建和分析查询三个阶段，
上下游之间需要传递数据集、文件、三元组和查询表等信息。
项目让每个阶段产生持久化标识，
下一个阶段凭标识读取上一个阶段的结构化结果：

\begin{itemize}
  \item 任务一向 DataMate 登记交付数据集，返回数据集标识；
  \item 任务二接收数据集标识，读取实际文件，
        入库时为每条三元组记录来源标识和文件标识；
  \item 任务三每次分析生成分析标识，
        将问题、SQL、数据行和图表说明打包为结构化对象。
\end{itemize}

这套标识串联形成三个具体结果：
第一，任务二不需要知道任务一用了哪些算子，
只需要一个数据集标识就能读取清洗结果；
第二，任务三的图表和导出文件均按同一个分析标识读取结果记录中的数据行；
第三，三元组的来源标识可用于查询原始文件，
分析标识可用于核对对应的查询计划和 SQL。

\section{离线优先的级联抽取}

医疗实体和关系抽取是任务二的核心环节。
将模型调用作为唯一抽取路径会带来两个实际限制：
一是模型服务波动或超时时整个抽取管道会中断；
二是批量处理会增加模型调用次数和等待时间，
在公开知识资源批量导入时成本很高。

MediFlow 为抽取服务提供三种执行方式：
本地词典与规则、大模型、以及离线优先的级联方式。
离线方式加载独立维护的约 3.17 万条医学术语词典和关系规则，
不调用外部模型，识别结果在相同输入和词典版本下可复现；
模型方式处理需要语义判断的文本；
级联方式先完成离线全量扫描，再按候选可靠性和句子覆盖情况决定模型复核或补充的范围。
三种方式通过同一组实体、关系和三元组字段提供结果，
上层编排逻辑不需要区分当前使用的是哪种方式。

在批量导入公开知识时可以显式使用离线方式，减少外部模型依赖；
任务二在线入口默认使用级联方式，
高可靠事实直接进入统一校验，中可靠候选进入复核，低可靠候选不进入主图谱；
对离线链路未覆盖但含有医疗提示词的句子，模型只返回原文直接支持的补充事实。
模型服务不可用时，系统保留已完成的离线结果并返回错误状态，
接口超时只影响对应的模型补充步骤。

在固定划分的离线评测中，CBLUE 的 CMeEE 实体识别评测集包含 2\,500 条样本，
实体识别精确率为 56.50\%、召回率为 29.57\%、F1 为 38.82\%；
CMeIE 评测集包含 1\,793 条样本，
关系抽取精确率为 11.46\%、召回率为 17.93\%、F1 为 13.98\%。
该结果用于说明当前离线链路的覆盖范围，
级联设计则把模型调用限制在缺口句和需复核候选上，
同时保留离线结果的可复现性、来源字段和原文片段。

\section{结构保持的数据处理}

项目输入同时包含 CSV 表格、JSON 对象和 JSONL 记录。
系统没有把所有输入先合并为无结构文本，
而是按格式保留列名、键名、嵌套层次和记录边界，
使后续知识抽取仍可按字段读取内容。

结构保持规则如下：
CSV 清洗时保留表头和行列位置，
JSON 清洗时保留键名和嵌套层次，
JSONL 清洗时保留逐行记录边界。
清洗完成后，汇聚步骤会对各格式产物重新解析检查：
CSV 重新读取表头确认列数一致，
JSON 和 JSONL 重新解析确认结构合法，
通过后才登记到最终数据集。
PDF 解析为文本后也保留源文件名，
知识抽取时可以区分哪些事实来自研文献、哪些来自表格数据。

由此，任务二在读取清洗结果时，
可以按字段取特定列的内容，
可以知道这条记录来自 CSV 的哪一行哪个文件。
这些信息会随来源字段和原文片段进入知识抽取结果。

\section{只读查询安全}

任务三需要根据自然语言问题执行 SQL 查询。
NL2SQL 的研究重点通常放在生成 SQL 的准确率上，
在一个需要公开展示的系统中，
SQL 的执行安全同样关键。
错误的 DELETE、多语句拼接、或者对系统表的访问，
都可能造成数据损坏或信息泄露。

MediFlow 将安全限制落实在三个层面：

\begin{itemize}
  \item \textbf{语句层}：只接受单条 SELECT 或 WITH 语句，
        拒绝多语句、注释注入和结构修改关键字；
  \item \textbf{对象层}：表名必须落在 17 张业务表和统计视图的允许清单中，
        跨库附加、系统表访问和自定义函数调用在解析阶段拦截；
  \item \textbf{连接层}：SQLite 以 URI 只读模式打开，
        连接后立即设置 \texttt{PRAGMA query\_only=ON}，
        即使前两层被绕过，数据库引擎本身也会拒绝写操作。
\end{itemize}

此外，每次分析最多执行 4 项查询、每项最多返回 200 行、
每项查询安装 5 秒进度中止条件。
这些限制的直接目的在于限制单次请求的资源占用，
避免异常查询长时间占用数据库连接。

这套三层安全方案同时服务于 Nexent 查询工具和独立工作台，
两个入口使用同一组执行规则。
八类业务查询的连续三轮测试（第七章）表明，
响应时间和返回行数见表~\ref{tab:online-query-results}。
